\section{气体动理论}
\sourceproblem{选择1}{54}{60}
\problemtext 刚性氢分子理想气体每升高1 K内能增加20.8 J，求气体质量。
\workings $\Delta U=(5/2)\nu R\Delta T$，故 $\nu=20.8/(2.5\times8.31)\approx1.00$ mol，质量约 $2.0\times10^{-3}$ kg，选B。
\crosscheck 采用氢气摩尔质量2 g/mol，忽略振动自由度。
\sourceproblem{选择2}{54}{60}
\problemtext 温度 $T$ 的刚性双原子氢气分子的平均动能。
\workings 三个平动、两个转动自由度，每个二次型自由度平均能量 $k_BT/2$，故平均总动能 $5k_BT/2$，选B。
\crosscheck 若只问平动动能则为 $3k_BT/2$；单分子用 $k_B$ 而非 $R$。
\sourceproblem{选择3}{54}{60}
\problemtext 氢气和氮气质量密度相同、平均平动动能相同且均平衡，比较温度和压强。
\workings $\overline K_{\rm tr}=3k_BT/2$，故同温；$p=\rho RT/M$，氮气摩尔质量较大，压强较小，选D。
\crosscheck 同质量密度并不意味着分子数密度相同。
\sourceproblem{选择4}{54}{60}
\problemtext 三容器分子数密度比4:2:1，平均平动动能比1:2:4，求压强比。
\workings $p=2n\overline K_{\rm tr}/3$，三者乘积相同，故压强比1:1:1，选A。
\crosscheck $n$ 为数密度，不是物质的量。
\sourceproblem{选择5}{54}{60}
\problemtext 判断温度与平动动能、分子热运动统计意义及单分子冷热的四个陈述。
\workings 平衡经典理想气体的温度量度平均平动动能，是大量分子的统计性质，并反映热运动程度；不能把热力学温度赋予每个瞬时分子。因此(1)(2)(3)正确，选B。
\crosscheck 单个分子动能随碰撞波动；平均动能与瞬时动能不同。
\sourceproblem{选择6}{54}{60}
\problemtext 1 mol单原子理想气体在1 atm下从0摄氏度升温到100摄氏度，求内能改变（选项续PDF61）。
\workings $\Delta U=3R\Delta T/2=1246.5$ J，约 $1.25\times10^3$ J，选D。
\crosscheck 等压过程的吸热量另含膨胀功，不能以 $C_p$ 计算内能变化。
\sourceproblem{选择7}{55}{61}
\problemtext 体积 $10^{-2}$ 立方米、气体质量100 g、方均根速率200 m/s，求压强。
\workings $\rho=10\,\mathrm{kg/m^3}$，$p=\rho v_{\rm rms}^2/3=1.33\times10^5$ Pa，选B。
\crosscheck 质量须先换算成千克。
\sourceproblem{选择8}{55}{61}
\problemtext 同温氧气、氢气麦克斯韦速率分布，较右峰速率2000 m/s，求氧气最概然速率。
\workings $v_p=\sqrt{2RT/M}$，右峰对应氢气，$v_{p,\mathrm O_2}=2000\sqrt{2/32}=500$ m/s，选A。
\crosscheck 同温平均平动动能相等而最概然速率不同。
\sourceproblem{选择9}{55}{61}
\problemtext 总分子数 $N$、归一化速率密度 $f(v)$，求速率大于100 m/s的分子数表达式。
\workings $N_{>100}=N\int_{100}^{\infty}f(v)\mathrm dv$，选B。
\crosscheck 不乘 $v$；积分本身是比例。
\sourceproblem{选择10}{55}{61}
\problemtext 定容固定量理想气体温度变为4倍，比较平均速率、平均碰撞频率和平均自由程。
\workings 硬球稀薄气体模型中 $\overline v\propto\sqrt T$，$\ell=1/(\sqrt2 n\pi d^2)$，$Z=\overline v/\ell$。数密度和碰撞直径不变，故速率、频率均为2倍，自由程不变，选B。
\crosscheck 若有效碰撞截面随温度变化，则自由程未必恒定；题目采用硬球模型。
